% \section{Related Work}
% \label{sec:relatedwork}
% Zhang et al.\cite{zhang1997birch} introduced the Clustering Feature Tree (CFT), known for its incrementality and additivity, making it ideal for streaming workloads. Aggarwal et al.~\cite{Clustream:03} extended the CF structure into microclusters (MCs), incorporating timestamps and weights to work with window models. Chen et al.~\cite{DStream:2007} proposed a grid-based data structure for efficiency. Gong et al.~\cite{EDMStream:17} presented the Dependency Tree (DPT), balancing efficiency and accuracy, although it may struggle with complex evolving features. 
% Various window models have been proposed to process subsets of data streams. Metwally et al.~\cite{metwally2005duplicate} introduced the landmark window model, while Zhou et al.~\cite{zhou2008tracking} and Borassi et al.~\cite{SL-KMeans:20} presented the sliding window model. Cao et al.~\cite{cao2006density} and Chen et al.~\cite{DStream:2007} proposed the damped window model, which retains all data but prioritizes recent information with varying weights. Outlier detection is crucial in \dsc algorithms. Early contributions like BIRCH~\cite{BIRCH:96} included a phase for identifying outliers based on density thresholds. Aggarwal et al.~\cite{Clustream:03} introduced the outlier timer for enhanced identification. Wan et al.~\cite{DenStream:06} conceptualized the outlier buffer to manage the transition between outliers and clustered points. Refinement strategies have been integral to \dsc algorithms since Aggarwal et al.~\cite{Clustream:03} introduced the online-offline clustering paradigm. Common offline clustering algorithms include k-means~\cite{KMeans:1967} and its variants such as scalable k-means~\cite{bradley2002scaling} and streaming k-means~\cite{farnstrom2000scalability}. Our empirical evaluations suggest that refinement strategies often introduce unnecessary computational overhead. These contributions inform \algo's automated selection of four design aspects of \dsc algorithms.


\section{Related Work} 
\label{sec:relatedwork} 
Data stream clustering algorithms have evolved significantly, yet many still lack the flexibility and self-optimization needed for diverse applications, resulting in suboptimal performance. Early work by Zhang et al.~\cite{zhang1997birch} introduced the Clustering Feature Tree (CFT) with a static structure, limiting adaptability. BIRCH~\cite{BIRCH:96} identified outliers based on density thresholds but did not adapt well to evolving data. Metwally et al.~\cite{metwally2005duplicate} proposed the landmark window model, while Zhou et al.~\cite{zhou2008tracking} and Borassi et al.~\cite{SL-KMeans:20} introduced variations of the sliding window model. These models improve adaptability but lack dynamic adjustment mechanisms. Aggarwal et al.~\cite{Clustream:03} extended CFT into microclusters (MCs) and introduced the online-offline clustering paradigm and the outlier timer, yet their fixed parameters limit adaptability. Cao et al.~\cite{cao2006density} and Chen et al.~\cite{DStream:2007} proposed the damped window model for temporal relevance, retaining all data but prioritizing recent information, though these models can be computationally expensive. Chen et al.~\cite{DStream:2007} also introduced a grid-based structure for efficiency. Wan et al.~\cite{DenStream:06} developed the outlier buffer, improving outlier management but adding computational overhead. Gong et al.~\cite{EDMStream:17} presented the Dependency Tree (DPT), balancing efficiency and accuracy but struggling with complex features. Farnstrom et al.~\cite{farnstrom2000scalability} and Bradley et al.~\cite{bradley2002scaling} proposed scalable and streaming k-means variants for large-scale data, yet these methods lack mechanisms for dynamic adjustment to data stream variations. These limitations highlight the need for a more flexible and adaptive approach. Our proposed algorithm, \algo, addresses these gaps by offering a modular and self-optimizing framework that dynamically balances clustering accuracy and efficiency, outperforming existing methods in varied data stream scenarios.







% \algo selects the most fitting window model based on user-defined objectives and data stream characteristics.

% \algo incorporates these mechanisms to choose the most suitable outlier detection strategy.

% \algo selects the most appropriate refinement strategy, aligned with user objectives and data stream characteristics.

 

% \textbf{Summarizing Data Structure.}
% The quest for efficient data structures for summarizing data streams has been a focal point in research. Zhang et al.\cite{zhang1997birch} introduced the Clustering Feature Tree (\CFT), characterized by its incrementality and additivity, making it well-suited for streaming workloads. Aggarwal et al.\cite{Clustream:03} extended the CF structure into microclusters (\MCs), incorporating additional summary information such as timestamps and weights to synergize with window models. Conversely, Chen et al.\cite{DStream:2007} advocated for a grid-based data structure for efficiency. Gong et al.\cite{EDMStream:17} presented the Dependency Tree (\DPT), which strikes a balance between efficiency and accuracy. However, \DPT may yield sub-optimal results in handling \cef effectively. These contributions inform \algo's automated selection of appropriate summarizing data structures.

% \textbf{Window Model.}
% Various window models have been proposed to specify the subset of data streams to be processed. Metwally et al.\cite{metwally2005duplicate} proposed the landmark window model (\landmark), while Zhou et al.\cite{zhou2008tracking} and Borassi et al.\cite{SL-KMeans:20} introduced the sliding window model (\sliding). Cao et al.\cite{cao2006density} and Chen et al.~\cite{DStream:2007} presented the damped window model (\damped), which retains all data but prioritizes the most recent information by associating varying weights. \algo leverages these foundational works to automatically select the most fitting window model, contingent on user-defined objectives and data stream characteristics.

% \textbf{Outlier Detection Mechanism.}
% Outlier detection is a pivotal design aspect in \dsc algorithms. Early contributions like \birch~\cite{BIRCH:96} included an optional phase for identifying outlier candidates based on object density thresholds. Aggarwal et al.\cite{Clustream:03} introduced the \emph{outlier timer} (\timer) for enhanced outlier identification. Wan et al.\cite{DenStream:06} further conceptualized the \emph{outlier buffer} (\buffer) to facilitate the transition between outliers and clustered points. \algo incorporates these mechanisms, enabling automated selection of the most suitable outlier detection strategy based on user objectives and data stream characteristics.

% \textbf{Refinement Strategy.}
% Refinement strategies have been integral to \dsc algorithms since Aggarwal et al.\cite{Clustream:03} introduced the online-offline clustering paradigm. Commonly employed offline clustering algorithms include \kmeans\cite{KMeans:1967} and its variants such as \sckmeans~\cite{bradley2002scaling} and \sinkmeans~\cite{farnstrom2000scalability}. However, our empirical evaluations suggest that refinement strategies often introduce unnecessary computational overhead. \algo takes these observations into account when automatically selecting the most appropriate refinement strategy, aligned with user objectives and data stream characteristics.